% ==========================================
% TIKZ FEYNMAN DIAGRAM DEFINITIONS
% ==========================================

% -----------------------------------------------------------------
% Wrapper function to scale down the diagrams (for inline use)
% -----------------------------------------------------------------

\newcommand{\scalefeyn}[2]{\scalebox{#1}{#2}}

% ---------------------------------------------------------
% CATEGORY 1: BASIC TOPOLOGIES & BUILDING BLOCKS (phi^4 theory)
% ---------------------------------------------------------

\newcommand{\feynLine}{%
    \tikz[baseline=-0.6ex, scale=0.5, thick]{
        \draw (0,0) -- (1.5,0);
        \filldraw (0,0) circle (2pt); \filldraw (1.5,0) circle (2pt);
    }%
}

\newcommand{\feynHoriz}{%
    \tikz[baseline=-0.6ex, scale=0.5, thick]{
        \draw (0,0.4) -- (1,0.4);
        \draw (0,-0.4) -- (1,-0.4);
        \filldraw (0,0.4) circle (2pt); \filldraw (1,0.4) circle (2pt);
        \filldraw (0,-0.4) circle (2pt); \filldraw (1,-0.4) circle (2pt);
    }%
}

\newcommand{\feynVert}{%
    \tikz[baseline=-0.6ex, scale=0.5, thick]{
        \draw (0.2,0.4) -- (0.2,-0.4);
        \draw (0.8,0.4) -- (0.8,-0.4);
        \filldraw (0.2,0.4) circle (2pt); \filldraw (0.2,-0.4) circle (2pt);
        \filldraw (0.8,0.4) circle (2pt); \filldraw (0.8,-0.4) circle (2pt);
    }%
}

\newcommand{\feynCross}{%
    \tikz[baseline=-0.6ex, scale=0.5, thick]{
        \draw (0,0.4) -- (1,-0.4);
        \draw (0,-0.4) -- (1,0.4);
        \filldraw (0,0.4) circle (2pt); \filldraw (1,-0.4) circle (2pt);
        \filldraw (0,-0.4) circle (2pt); \filldraw (1,0.4) circle (2pt);
    }%
}

% ---------------------------------------------------------
% CATEGORY 2: 4-POINT SCATTERING DIAGRAMS (phi^4 theory)
% ---------------------------------------------------------

% 0. 4-point tree (single vertex)
\newcommand{\feynFourTree}{%
    \tikz[baseline=-0.6ex, scale=0.6, thick]{
        % External legs
        \draw (0,0) -- (-1.2, 0.5) node[left] {\(x_1\)};
        \draw (0,0) -- (-1.2,-0.5) node[left] {\(x_2\)};
        \draw (0,0) -- (1.2, 0.5) node[right] {\(x_3\)};
        \draw (0,0) -- (1.2,-0.5) node[right] {\(x_4\)};
        % Vertex
        \filldraw (0,0) circle (1.5pt);
    }%
}

% 1. 4-point Bubble (Evolution of feynLoop)
\newcommand{\feynFourBubble}{%
    \tikz[baseline=-0.6ex, scale=0.6, thick]{
        \draw (0,0) circle (0.5); % Central loop
        % External legs
        \draw (-0.5,0) -- (-1.2, 0.5) node[left] {\(x_1\)};
        \draw (-0.5,0) -- (-1.2,-0.5) node[left] {\(x_2\)};
        \draw (0.5,0) -- (1.2, 0.5) node[right] {\(x_3\)};
        \draw (0.5,0) -- (1.2,-0.5) node[right] {\(x_4\)};
        % Dots on internal vertices
        \filldraw (-0.5,0) circle (1.5pt);
        \filldraw (0.5,0) circle (1.5pt);
        % Dots on external points
        \filldraw (-1.2,0.5) circle (1.5pt); \filldraw (-1.2,-0.5) circle (1.5pt);
        \filldraw (1.2,0.5) circle (1.5pt); \filldraw (1.2,-0.5) circle (1.5pt);
    }%
}

% 2. 4-point Sunset (Sunset insertion on s-channel)
\newcommand{\feynFourSunset}{%
    \tikz[baseline=-0.6ex, scale=0.6, thick]{
        % Central Sunset insertion
        \draw (0,0) circle (0.5);
        \draw (-0.5,0) -- (0.5,0);
        % Internal propagators to scattering vertices
        \draw (-0.5,0) -- (-1.5,0);
        \draw (0.5,0) -- (1.5,0);
        % External legs
        \draw (-1.5,0) -- (-2.2, 0.5) node[left] {\(x_1\)};
        \draw (-1.5,0) -- (-2.2,-0.5) node[left] {\(x_2\)};
        \draw (1.5,0) -- (2.2, 0.5) node[right] {\(x_3\)};
        \draw (1.5,0) -- (2.2,-0.5) node[right] {\(x_4\)};
        % Dots
        \filldraw (-0.5,0) circle (1.5pt); \filldraw (0.5,0) circle (1.5pt);
        \filldraw (-1.5,0) circle (1.5pt); \filldraw (1.5,0) circle (1.5pt);
        \filldraw (-2.2,0.5) circle (1.5pt); \filldraw (-2.2,-0.5) circle (1.5pt);
        \filldraw (2.2,0.5) circle (1.5pt); \filldraw (2.2,-0.5) circle (1.5pt);
    }%
}

% 3. 4-point Dumbbell (Double Tadpole insertion on s-channel)
\newcommand{\feynFourDumbbell}{%
    \tikz[baseline=-0.6ex, scale=0.6, thick]{
        % Central line where loops rest (the dumbbell bar)
        \draw (-1.5,0) -- (1.5,0);
        % The two loops (tadpoles)
        \draw (-0.5,0.4) circle (0.4);
        \draw (0.5,0.4) circle (0.4);
        % External legs
        \draw (-1.5,0) -- (-2.2, 0.5) node[left] {\(x_1\)};
        \draw (-1.5,0) -- (-2.2,-0.5) node[left] {\(x_2\)};
        \draw (1.5,0) -- (2.2, 0.5) node[right] {\(x_3\)};
        \draw (1.5,0) -- (2.2,-0.5) node[right] {\(x_4\)};
        % Dots
        \filldraw (-0.5,0) circle (1.5pt); \filldraw (0.5,0) circle (1.5pt);
        \filldraw (-1.5,0) circle (1.5pt); \filldraw (1.5,0) circle (1.5pt);
        \filldraw (-2.2,0.5) circle (1.5pt); \filldraw (-2.2,-0.5) circle (1.5pt);
        \filldraw (2.2,0.5) circle (1.5pt); \filldraw (2.2,-0.5) circle (1.5pt);
    }%
}

% ---------------------------------------------------------
% CATEGORY 4: SPECIFIC 2-POINT CORRELATION FUNCTIONS
% ---------------------------------------------------------

% 1. Tree Level (Free propagator)
\newcommand{\feynTree}{%
    \tikz[baseline=-0.6ex, scale=0.5, thick]{
        \draw (0,0) -- (2,0);
        \filldraw (0,0) circle (2pt) node[below=2pt] {\(x_1\)};
        \filldraw (2,0) circle (2pt) node[below=2pt] {\(x_2\)};
    }%
}

% 2. 1-Loop Bubble
\newcommand{\feynBubble}{%
    \tikz[baseline=-0.6ex, scale=0.5, thick]{
        \draw (0,0) -- (0.8,0); % left leg
        \draw (1.4,0) circle (0.6); % central bubble (center 1.4, radius 0.6)
        \draw (2.0,0) -- (2.8,0); % right leg

        \filldraw (0,0) circle (2pt) node[below=2pt] {\(x_1\)};
        \filldraw (0.8,0) circle (2pt);
        \filldraw (2.0,0) circle (2pt);
        \filldraw (2.8,0) circle (2pt) node[below=2pt] {\(x_2\)};
    }%
}

% 3. Double Tadpole (two tadpoles on the propagator)
\newcommand{\feynDoubleTadpole}{%
    \tikz[baseline=-0.6ex, scale=0.5, thick]{
        \draw (0,0) -- (0.8,0); % first external propagator
        \draw (2,0) -- (2.8,0); % second external propagator
        \draw (0.8,0.4) circle (0.4); % first loop
        \draw (2.0,0.4) circle (0.4); % second loop

        \filldraw (0,0) circle (2pt) node[below=2pt] {\(x_1\)};
        \filldraw (0.8,0) circle (2pt);
        \filldraw (2.0,0) circle (2pt);
        \filldraw (2.8,0) circle (2pt) node[below=2pt] {\(x_2\)};
    }%
}

% 4. Cactus (Tadpole on top of another Tadpole)
\newcommand{\feynCactus}{%
    \tikz[baseline=-0.6ex, scale=0.5, thick]{
        \draw (0,0) -- (2.4,0); % horizontal line
        \draw (1.2,0) -- (1.2,0.8); % cactus stem
        \draw (1.2,1.2) circle (0.4); % top loop

        \filldraw (0,0) circle (2pt) node[below=2pt] {\(x_1\)};
        \filldraw (1.2,0) circle (2pt);
        \filldraw (1.2,0.8) circle (2pt);
        \filldraw (2.4,0) circle (2pt) node[below=2pt] {\(x_2\)};
    }%
}

% 5. Disconnected 1: Propagator + Double Tadpole in the vacuum
\newcommand{\feynDiscDoubleTadpole}{%
    \tikz[baseline=-0.6ex, scale=0.5, thick]{
        % Connected part at the bottom
        \draw (0,0) -- (2.4,0);
        \filldraw (0,0) circle (2pt) node[below=2pt] {\(x_1\)};
        \filldraw (2.4,0) circle (2pt) node[below=2pt] {\(x_2\)};

        % Disconnected part at the top
        \draw (0.7,1) -- (1.7,1); % horizontal stem
        \draw (0.7,1.3) circle (0.3); % left loop
        \draw (1.7,1.3) circle (0.3); % right loop
        \filldraw (0.7,1) circle (2pt);
        \filldraw (1.7,1) circle (2pt);
    }%
}

% 6. Disconnected 2: Propagator + Sunset (Melon) in the vacuum
\newcommand{\feynDiscSunset}{%
    \tikz[baseline=-0.6ex, scale=0.5, thick]{
        % Connected part at the bottom
        \draw (0,0) -- (2.4,0);
        \filldraw (0,0) circle (2pt) node[below=2pt] {\(x_1\)};
        \filldraw (2.4,0) circle (2pt) node[below=2pt] {\(x_2\)};

        % Disconnected part at the top (Sunset)
        \draw (1.2,1.1) circle (0.5); % outer circle
        \draw (0.7,1.1) -- (1.7,1.1); % central line
        \filldraw (0.7,1.1) circle (2pt);
        \filldraw (1.7,1.1) circle (2pt);
    }%
}

% ---------------------------------------------------------
% CATEGORY 5: 4-POINT CORRELATION FUNCTIONS
% ---------------------------------------------------------

% 1. s-channel (Horizontal propagator)
\newcommand{\feynFourS}{%
    \tikz[baseline=-0.6ex, scale=0.5, thick]{
        % Internal propagator
        \draw (0.8,0) -- (1.6,0);

        % Left legs
        \draw (0,0.8) -- (0.8,0);
        \draw (0,-0.8) -- (0.8,0);
        % Right legs
        \draw (1.6,0) -- (2.4,0.8);
        \draw (1.6,0) -- (2.4,-0.8);

        % Internal vertices
        \filldraw (0.8,0) circle (2pt);
        \filldraw (1.6,0) circle (2pt);
        % External nodes with labels
        \filldraw (0,0.8) circle (2pt) node[left=2pt] {\(x_1\)};
        \filldraw (0,-0.8) circle (2pt) node[left=2pt] {\(x_2\)};
        \filldraw (2.4,0.8) circle (2pt) node[right=2pt] {\(x_3\)};
        \filldraw (2.4,-0.8) circle (2pt) node[right=2pt] {\(x_4\)};
    }%
}

% 2. t-channel (Vertical propagator)
\newcommand{\feynFourT}{%
    \tikz[baseline=-0.6ex, scale=0.5, thick]{
        % Internal propagator
        \draw (1.2,0.4) -- (1.2,-0.4);

        % Upper legs
        \draw (0,0.8) -- (1.2,0.4);
        \draw (2.4,0.8) -- (1.2,0.4);
        % Lower legs
        \draw (0,-0.8) -- (1.2,-0.4);
        \draw (2.4,-0.8) -- (1.2,-0.4);

        % Internal vertices
        \filldraw (1.2,0.4) circle (2pt);
        \filldraw (1.2,-0.4) circle (2pt);
        % External nodes with labels
        \filldraw (0,0.8) circle (2pt) node[left=2pt] {\(x_1\)};
        \filldraw (0,-0.8) circle (2pt) node[left=2pt] {\(x_2\)};
        \filldraw (2.4,0.8) circle (2pt) node[right=2pt] {\(x_3\)};
        \filldraw (2.4,-0.8) circle (2pt) node[right=2pt] {\(x_4\)};
    }%
}

% 3. u-channel (Crossed legs)
\newcommand{\feynFourU}{%
    \tikz[baseline=-0.6ex, scale=0.5, thick]{
        % Left legs
        \draw (0,0.8) -- (1.2,0.4);
        \draw (0,-0.8) -- (1.2,-0.4);

        % Leg from bottom vertex to x3 (top right)
        \draw (1.2,-0.4) -- (2.4,0.8);

        % Trick: thick white line to "break" the underlying line
        \draw[white, line width=4pt] (1.2,0.4) -- (2.4,-0.8);
        % Leg from top vertex to x4 (bottom right)
        \draw (1.2,0.4) -- (2.4,-0.8);

        % Internal vertices
        \filldraw (1.2,0.4) circle (2pt);
        \filldraw (1.2,-0.4) circle (2pt);

        % Internal propagator
        \draw (1.2,0.4) -- (1.2,-0.4);

        % External nodes with labels
        \filldraw (0,0.8) circle (2pt) node[left=2pt] {\(x_1\)};
        \filldraw (0,-0.8) circle (2pt) node[left=2pt] {\(x_2\)};
        \filldraw (2.4,0.8) circle (2pt) node[right=2pt] {\(x_3\)};
        \filldraw (2.4,-0.8) circle (2pt) node[right=2pt] {\(x_4\)};
    }%
}

% 4. Disconnected: Vertical (x1-x2 and x3-x4)
\newcommand{\feynFourDiscVert}{%
    \tikz[baseline=-0.6ex, scale=0.5, thick]{
        % Free propagators
        \draw (0,0.8) -- (0,-0.8);
        \draw (2.4,0.8) -- (2.4,-0.8);

        % External nodes with labels
        \filldraw (0,0.8) circle (2pt) node[left=2pt] {\(x_1\)};
        \filldraw (0,-0.8) circle (2pt) node[left=2pt] {\(x_2\)};
        \filldraw (2.4,0.8) circle (2pt) node[right=2pt] {\(x_3\)};
        \filldraw (2.4,-0.8) circle (2pt) node[right=2pt] {\(x_4\)};
    }%
}

% 5. Disconnected: Horizontal (x1-x3 and x2-x4)
\newcommand{\feynFourDiscHoriz}{%
    \tikz[baseline=-0.6ex, scale=0.5, thick]{
        % Free propagators
        \draw (0,0.8) -- (2.4,0.8);
        \draw (0,-0.8) -- (2.4,-0.8);

        % External nodes with labels
        \filldraw (0,0.8) circle (2pt) node[left=2pt] {\(x_1\)};
        \filldraw (0,-0.8) circle (2pt) node[left=2pt] {\(x_2\)};
        \filldraw (2.4,0.8) circle (2pt) node[right=2pt] {\(x_3\)};
        \filldraw (2.4,-0.8) circle (2pt) node[right=2pt] {\(x_4\)};
    }%
}

% 6. Disconnected: Crossed (x1-x4 and x2-x3)
\newcommand{\feynFourDiscCross}{%
    \tikz[baseline=-0.6ex, scale=0.5, thick]{
        % First free propagator (x2 to x3)
        \draw (0,-0.8) -- (2.4,0.8);

        % Trick: thick white line to break the crossing
        \draw[white, line width=4pt] (0,0.8) -- (2.4,-0.8);

        % Second free propagator (x1 to x4)
        \draw (0,0.8) -- (2.4,-0.8);

        % External nodes with labels
        \filldraw (0,0.8) circle (2pt) node[left=2pt] {\(x_1\)};
        \filldraw (0,-0.8) circle (2pt) node[left=2pt] {\(x_2\)};
        \filldraw (2.4,0.8) circle (2pt) node[right=2pt] {\(x_3\)};
        \filldraw (2.4,-0.8) circle (2pt) node[right=2pt] {\(x_4\)};
    }%
}

% ---------------------------------------------------------
% CATEGORY 6: Amputated diagrams
% ---------------------------------------------------------

% 1. Sum of 2-point diagrams (Hatched bubble on a line)
\newcommand{\legFullProp}{%
    \vcenter{\hbox{\tikz[scale=0.8, thick]{
                % Draw the continuous line first
                \draw (0,0) -- (2,0);

                % Fill a white circle to hide the line behind the bubble
                \path[fill=white] (1,0) circle (0.4cm);

                % Draw the hatched bubble on top
                \node[circle, draw, thick, minimum size=0.7cm, pattern=north east lines, pattern color=black!80] at (1,0) {};
            }}}%
}

% 2. Sum of amputated diagrams (Blob with external legs)
\newcommand{\legAmputated}{%
    \vcenter{\hbox{\tikz[scale=0.8, thick]{
                % The central blob (Ellipse)
                \node[draw, thick, ellipse, fill=white, minimum width=1.6cm, minimum height=0.8cm, font=\small] (blob) at (0,0) {AMPUTATED};

                % Left legs (connecting to specific angles of the ellipse)
                \draw (-2.2, 0.6) -- (blob.160);
                \draw (-2.2, -0.6) -- (blob.200);

                % Right legs
                \draw (blob.20) -- (2.2, 0.6);
                \draw (blob.340) -- (2.2, -0.6);
            }}}%
}

% ---------------------------------------------------------
% CATEGORY 7: DYSON SERIES & 1PI BLOBS
% ---------------------------------------------------------

% 1. Empty line (Free propagator, no corrections)
\newcommand{\feynBare}{%
    \tikz[baseline=-0.6ex, scale=0.5, thick]{
        \draw (0,0) -- (2,0);
    }%
}

% 2. Bubble with lines (Generic 1PI insertion)
\newcommand{\feynBlobNE}{%
    \tikz[baseline=-0.6ex, scale=0.5, thick]{
        \draw (0,0) -- (3,0);
        \draw[fill=white, postaction={pattern=north east lines, pattern}] (1.5,0) circle (0.6);
    }%
}

% 3. Empty bubble (Single loop with one particle)
\newcommand{\feynBlobEmpty}{%
    \tikz[baseline=-0.6ex, scale=0.5, thick]{
        \draw (0,0) -- (3,0);
        \filldraw[fill=white, thick] (1.5,0) circle (0.6);
    }%
}

% 4. Bubble with vertical line
\newcommand{\feynBlobVert}{%
    \tikz[baseline=-0.6ex, scale=0.5, thick]{
        \draw (0,0) -- (2.5,0);
        \filldraw[fill=white, thick] (1.25,0) circle (0.6);
        \draw (1.25,-0.6) -- (1.25,0.6);
    }%
}

% ---------------------------------------------------------
% TEMPORARY: INLINE FEYNMAN PARSER (expl3 based)
% ---------------------------------------------------------

\ExplSyntaxOn
\NewDocumentCommand{\inlinefeynman}{m}
{
    % baseline aggiusta l'altezza del diagramma rispetto al testo
    \tikz[baseline=-0.5ex, x=0.6cm, y=0.6cm, line~width=0.6pt] {

        % l_tmpa_int è la coordinata X di partenza (inizializzata a 0)
        % l_tmpb_int è la coordinata X di arrivo del blocco corrente
        \int_zero:N \l_tmpa_int
        \int_zero:N \l_tmpb_int

        \tl_map_inline:nn { #1 }
        {
            % Ad ogni passo, calcoliamo la coordinata di arrivo (X + 1)
            \int_set:Nn \l_tmpb_int { \l_tmpa_int + 1 }

            \str_case:nn { ##1 }
            {
                { - } {
                        % Propagatore: linea dritta. Nodi a sinistra e a destra.
                        \fill (\int_use:N \l_tmpa_int, 0) circle (1.5pt);
                        \fill (\int_use:N \l_tmpb_int, 0) circle (1.5pt);
                        \draw (\int_use:N \l_tmpa_int, 0) -- (\int_use:N \l_tmpb_int, 0);
                    }
                    { o } {
                        % Anello (Bolla): due archi. Nodi a sinistra e a destra.
                        \fill (\int_use:N \l_tmpa_int, 0) circle (1.5pt);
                        \fill (\int_use:N \l_tmpb_int, 0) circle (1.5pt);
                        \draw (\int_use:N \l_tmpa_int, 0) to[out=90, in=90] (\int_use:N \l_tmpb_int, 0);
                        \draw (\int_use:N \l_tmpa_int, 0) to[out=-90, in=-90] (\int_use:N \l_tmpb_int, 0);
                    }
                    { > } {
                        % Linee entranti (convergono nel nodo a destra)
                        \fill (\int_use:N \l_tmpb_int, 0) circle (1.5pt);
                        \draw (\int_use:N \l_tmpa_int, 0.5) -- (\int_use:N \l_tmpb_int, 0);
                        \draw (\int_use:N \l_tmpa_int, -0.5) -- (\int_use:N \l_tmpb_int, 0);
                    }
                    { < } {
                        % Linee uscenti (divergono dal nodo a sinistra)
                        \fill (\int_use:N \l_tmpa_int, 0) circle (1.5pt);
                        \draw (\int_use:N \l_tmpa_int, 0) -- (\int_use:N \l_tmpb_int, 0.5);
                        \draw (\int_use:N \l_tmpa_int, 0) -- (\int_use:N \l_tmpb_int, -0.5);
                    }
            }
            % Aggiorniamo la coordinata di partenza per il prossimo elemento
            \int_set_eq:NN \l_tmpa_int \l_tmpb_int
        }
    }
}
\ExplSyntaxOff